\documentclass{article}
\usepackage[letterpaper,top=2.54cm,bottom=2.54cm,left=2.54cm,right=2.54cm,marginparwidth=1.75cm]{geometry}
\usepackage{url}
\usepackage{biblatex}
\addbibresource{bibliography.bib}
\begin{document}
\noindent GIF Research and Development Outlook for Gen IV, 2018

\noindent \cite{GIF2018Roadmap}
\begin{itemize}
    \item objectives are: 
    \begin{itemize}
        \item sustainability
        \item safety and reliability
        \item economic competitiveness
        \item proliferation resistance
        \item physical protection
    \end{itemize}
    \item most promising designs:
    \begin{itemize}
        \item sodium-cooled fast reactor (SFR)
        \item very high temperature reactor (VHTR)
        \item gas-cooled fast reactor (GFR)
        \item molten salt reactor (MSR)
        \item lead-cooled fast reactor (LFR)
        \item supercritical water-cooled reactor (SCWR)
    \end{itemize}
    \item viability, performance, demonstration/deployment
    \item high temp systems (VHTR, GFR, LFR and MSR) have potential applications hydrogen production or industrial process heat for such chemical processing facilities as petroleum refineries. 
    \item The three fast reactor systems and the MSR have the capability to transmute actinides for effective waste
management
\end{itemize}

Waste treatment that also has some types of waste
\url{https://www.sciencedirect.com/science/article/pii/S002954931930024X}
\begin{itemize}
    \item particulates, aerosols, reactive gases, tritium, residual halides (F, Cl, I), O2 \& N2, noble gases (Kr, Xe)
\end{itemize}


\section{Molten Salt Reactors}
Going down the thingamibob, e) says one of the challenges is on-site fuel processing. Specifically, interest in neutron poison and protactinium removal from carrier salt:

\begin{quote}

The present effort in electrochemistry is focused on the
development and verification of quantitative electrochemical extraction of uranium and
thorium on removal of main neutron poisons (fission products) from the MSR carrier salt
(LiF-BeF2). Special attention will be paid to the electrochemical studies of protactinium \cite{ReprocessingReductiveExtraction}.

\end{quote}

We want to remove Xenon and Samarium from the salt. Two-step thingymabob, first xenon removal \cite{XeRemoval} (prior to heat exchanger) and second is samarium removal \cite{SmRemoval} (before / after heat exchanger depending on conditions needed)

Review of Hazards Associated with Molten Salt Reactor Fuel Processing Operations (general)
https://www.nrc.gov/docs/ML1920/ML19205A386.pdf

149Sm evolution behavior in a small modular molten salt reactor
https://www.sciencedirect.com/science/article/pii/S0306454918302743

Electrochemical extraction of samarium from molten chlorides in pyrochemical processes
https://www.sciencedirect.com/science/article/pii/S0013468611011030
(USED THIS ONE)

Protact removal (dormant)
https://www.osti.gov/servlets/purl/1524586#:~:text=To%20obtain%20a%20high%20breeding,a%20higher%20protac%2D%20tinium%20concentration.

Helium bubbling in a Molten Salt Fast Reactor (???)
https://filelist.tudelft.nl/TNW/Afdelingen/Radiation%20Science%20and%20Technology/Research%20Groups/RPNM/Publications/MSc_Dirkjan_Journee.pdf

The chemistry and thermodynamics of molten salt reactor fuels (general)
https://www.sciencedirect.com/science/article/abs/pii/002231157490124X

Molten salt reactor fuels (general)
https://www.sciencedirect.com/science/article/abs/pii/B9780081025710000070


\printbibliography
\end{document}